\ifdefined\gkver\else\def\gkver{student}\fi
\documentclass[\gkver]{gaokaozhenti}
\begin{document}

\chapter{1998年上海}

\begin{enumerate}
\item
%% number: 1
%% paperName: 1998年上海高考第 1 题 4 分
%% typeId: 1
%% type: 单选题
%% chapter: 光学
%% point: 光电效应
%% method: 无
%% score: 4
%% degree: 850
%% duplicateId: 0
%% body:
下列实验中，能证实光具有粒子性的是\xzanswer{A}
\fourchoices[answer=A]
{光电效应实验}
{光的双缝干涉实验}
{光的圆孔衍射实验}
{$\alpha$ 粒子散射实验}

%% answer: A
%% memo:
\memoanswer{
本题原卷及材料中未提供详解，待补充。
}

\item
%% number: 2
%% paperName: 1998年上海高考第 2 题 4 分
%% typeId: 1
%% type: 单选题
%% chapter: 原子和原子核
%% point: 人工核反应
%% method: 无
%% score: 4
%% degree: 850
%% duplicateId: 0
%% body:
下列核反应方程中正确的是\xzanswer{B}
\fourchoices[answer=B]
{$\ce{^{238}_{92}U -> ^{234}_{90}Th + ^{2}_{1}H}$}
{$\ce{^{9}_{4}Be + ^{4}_{2}He -> ^{12}_{6}C + ^{1}_{0}n}$}
{$\ce{^{234}_{90}Th -> ^{234}_{90}Pa + ^{0}_{-1}e}$}
{$\ce{^{31}_{15}P -> ^{30}_{14}Si + ^{0}_{1}e}$}

%% answer: B
%% memo:
\memoanswer{
本题原卷及材料中未提供详解，待补充。
}

\item
%% number: 3
%% paperName: 1998年上海高考第 3 题 4 分
%% typeId: 1
%% type: 单选题
%% chapter: 直线运动
%% point: 运动图象
%% method: 无
%% score: 4
%% degree: 650
%% duplicateId: 0
%% body:
有两个光滑固定的斜面 AB 和 BC，A 和 C 两点在同一水平面上，斜面 BC 比斜面 AB 长。一个滑块自 A 点以速度 $v_{A}$ 上滑，到达 B 点时速度减小为零，紧接着沿 BC 滑下。设滑块从 A 点到 C 点的总时间是 $t_{C}$，那么下列四个图中，正确表示滑块速度的大小 $v$ 随时间 $t$ 变化的规律的是\xzanswer{C}

\onepicture[w=5cm]{\includesvg{figs/03a}}

\fourchoices[answer=C, ispicture=true, h=2.2cm]
{figs/03b.pdf}
{figs/03c.pdf}
{figs/03d.pdf}
{figs/03e.pdf}

%% answer: C
%% memo:
\memoanswer{
本题原卷及材料中未提供详解，待补充。
}

\item
%% number: 4
%% paperName: 1998年上海高考第 4 题 4 分
%% typeId: 1
%% type: 单选题
%% chapter: 电路
%% point: 闭合电路欧姆定律
%% method: 无
%% score: 4
%% degree: 650
%% duplicateId: 0
%% body:
如图所示电路中，电源 $E$ 的电动势为 $3.2\UV$，电阻 $R$ 的阻值为 $30\UO$，小灯泡 L 的额定电压为 $3.0\UV$，额定功率为 $4.5\UW$。当电键 S 接位置 1 时，电压表的读数为 $3\UV$，那么当电键 S 接到位置 2 时，小灯泡 L 的发光情况是\xzanswer{A}

\onepicture[w=6cm]{\includesvg{figs/04}}

\fourchoices[answer=A]
{很暗，甚至不亮}
{正常发光}
{比正常发光略亮}
{有可能被烧坏}

%% answer: A
%% memo:
\memoanswer{
本题原卷及材料中未提供详解，待补充。
}

\item
%% number: 5
%% paperName: 1998年上海高考第 5 题 4 分
%% typeId: 1
%% type: 单选题
%% chapter: 交流电
%% point: LC振荡回路
%% method: 无
%% score: 4
%% degree: 650
%% duplicateId: 0
%% body:
如图所示电路中，L 是电阻不计的电感器，C 是电容器。闭合电键 S，待电路达到稳定状态后，再打开电键 S，LC 电路中将产生电磁振荡。如果规定电感 L 中的电流方向从 a 到 b 为正，打开电键的时刻 $t = 0$，那么下列四个图中能正确表示电感中的电流 $i$ 随时间 $t$ 变化规律的是\xzanswer{B}

\onepicture[w=4.5cm]{\includesvg{figs/05a}}

\fourchoices[answer=B, ispicture=true, h=2.2cm]
{figs/05b.pdf}
{figs/05c.pdf}
{figs/05d.pdf}
{figs/05e.pdf}

%% answer: B
%% memo:
\memoanswer{
本题原卷及材料中未提供详解，待补充。
}

\item
%% number: 6
%% paperName: 1998年上海高考第 6 题 4 分
%% typeId: 1
%% type: 单选题
%% chapter: 力物体的平衡
%% point: 物体系平衡
%% method: 整体与隔离
%% score: 4
%% degree: 650
%% duplicateId: 0
%% body:
有一个直角支架 AOB，AO 水平放置，表面粗糙，OB 竖直向下，表面光滑。AO 上套有小环 P，OB 套有小环 Q，两环质量均为 $m$，两环间由一根质量可忽略、不可伸长的细绳相连，并在某一位置平衡，如图所示。现将 P 环向左移一小段距离，两环再次达到平衡，那么将 P 环移动后的平衡状态和原来的平衡状态比较，AO 杆对 P 环的支持力 $N$ 和细绳上拉力 $T$ 的变化情况是\xzanswer{B}

\onepicture[w=4.5cm]{\includesvg{figs/06}}

\fourchoices[answer=B]
{$N$ 不变，$T$ 变大}
{$N$ 不变，$T$ 变小}
{$N$ 变大，$T$ 变大}
{$N$ 变大，$T$ 变小}

%% answer: B
%% memo:
\memoanswer{
本题原卷及材料中未提供详解，待补充。
}

\item
%% number: 7
%% paperName: 1998年上海高考第 7 题 5 分
%% typeId: 2
%% type: 多选题
%% chapter: 气体性质
%% point: 热力学第一定律
%% method: 无
%% score: 5
%% degree: 650
%% duplicateId: 0
%% body:
有关物体内能，以下说法正确的是\xzanswer{AD}
\fourchoices[answer=AD]
{$1\Ug$ $0\UdC$ 水的内能比 $1\Ug$ $0\UdC$ 冰的内能大}
{电流通过电阻后发热，它的内能增加是通过“热传递”方式实现的}
{气体膨胀，它的内能一定减少}
{橡皮筋被拉伸时，分子间势能增加}

%% answer: AD
%% memo:
\memoanswer{
本题原卷及材料中未提供详解，待补充。
}

\item
%% number: 8
%% paperName: 1998年上海高考第 8 题 5 分
%% typeId: 2
%% type: 多选题
%% chapter: 运动和力
%% point: 动量定理
%% method: 无
%% score: 5
%% degree: 650
%% duplicateId: 0
%% body:
在光滑水平面上有质量均为 $2\Ukg$ 的 a、b 两质点，a 质点在水平恒力 $F_{a} = 4\UN$ 的作用下由静止出发运动 $4\Us$，b 质点在水平恒力 $F_{b} = 4\UN$ 作用下由静止出发移动 $4\Um$。比较这两质点所经历的过程，可以得到的结论是\xzanswer{AC}
\fourchoices[answer=AC]
{a 质点的位移比 b 质点的位移大}
{a 质点的末速度比 b 质点的末速度小}
{力 $F_{a}$ 作的功比力 $F_{b}$ 作的功多}
{力 $F_{a}$ 的冲量比力 $F_{b}$ 的冲量小}

%% answer: AC
%% memo:
\memoanswer{
本题原卷及材料中未提供详解，待补充。
}

\item
%% number: 9
%% paperName: 1998年上海高考第 9 题 5 分
%% typeId: 2
%% type: 多选题
%% chapter: 电磁感应
%% point: 楞次定律
%% method: 无
%% score: 5
%% degree: 450
%% duplicateId: 0
%% body:
如图所示，在一固定圆柱形磁铁的 N 极附近置一平面线圈 abcd，磁铁轴线与线圈水平中心线 $xx'$ 轴重合。下列说法正确的是\xzanswer{CD}

\onepicture[w=5.5cm]{\includesvg{figs/09}}

\fourchoices[answer=CD]
{当线圈刚沿 $xx'$ 轴向右平移时，线圈中有感应电流}
{当线圈刚绕 $xx'$ 轴转动时（ad 向外，bc 向里），线圈中有感应电流}
{当线圈刚沿垂直纸面方向向外平移时，线圈中有感应电流}
{当线圈刚绕 $yy'$ 轴转动时（ab 向里，cd 向外），线圈中有感应电流}

%% answer: CD
%% memo:
\memoanswer{
本题原卷及材料中未提供详解，待补充。
}

\item
%% number: 10
%% paperName: 1998年上海高考第 10 题 5 分
%% typeId: 2
%% type: 多选题
%% chapter: 电场
%% point: 带电粒子的轨迹类问题
%% method: 无
%% score: 5
%% degree: 450
%% duplicateId: 0
%% body:
图中实线是一簇未标明方向的由点电荷产生的电场线，虚线是某一带电粒子通过该电场区域时的运动轨迹，a、b 是轨迹上的两点。若带电粒子在运动中只受电场力作用，根据此图可作出正确判断的是\xzanswer{BCD}

\onepicture[w=4.5cm]{\includesvg{figs/10}}

\fourchoices[answer=BCD]
{带电粒子所带电荷的符号}
{带电粒子在 a、b 两点的受力方向}
{带电粒子在 a、b 两点的速度何处较大}
{带电粒子在 a、b 两点的电势能何处较大}

%% answer: BCD
%% memo:
\memoanswer{
本题原卷及材料中未提供详解，待补充。
}

\item
%% number: 11
%% paperName: 1998年上海高考第 11 题 5 分
%% typeId: 2
%% type: 多选题
%% chapter: 曲线运动
%% point: 人造地球卫星
%% method: 无
%% score: 5
%% degree: 450
%% duplicateId: 0
%% body:
发射地球同步卫星时，先将卫星发射至近地圆轨道 1，然后点火，使其沿椭圆轨道 2 运行，最后再次点火。将卫星送入同步圆轨道 3。轨道 1、2 相切于 Q 点，轨道 2、3 相切于 P 点（如图），则当卫星分别在 1，2，3 轨道上正常运行时，以下说法正确的是\xzanswer{BD}

\onepicture[w=5cm]{\includesvg{figs/11}}

\fourchoices[answer=BD]
{卫星在轨道 3 上的速率大于在轨道 1 的速率}
{卫星在轨道 3 上的角速度小于在轨道 1 上的角速度}
{卫星在轨道 1 上经过 Q 点时的加速度大于它在轨道 2 上经过 Q 点时的加速度}
{卫星在轨道 2 上经过 P 点的加速度等于它在轨道 3 上经过 P 点的加速度}

%% answer: BD
%% memo:
\memoanswer{
本题原卷及材料中未提供详解，待补充。
}

\item
%% number: 12
%% paperName: 1998年上海高考第 12 题 4 分
%% typeId: 3
%% type: 填空题
%% chapter: 原子和原子核
%% point: 人工核反应
%% method: 无
%% score: 4
%% degree: 850
%% duplicateId: 0
%% body:
现在，科学家们正在设法探寻“反物质”。所谓“反物质”是由“反粒子”构成的。“反粒子”与其对应的正粒子具有相同的质量和相同的电量，但电荷的符号相反。据此，若有反 $\alpha$ 粒子，它的质量数为 \tkanswer[1.2cm]{4}，电荷数为 \tkanswer[1.2cm]{$-2$}。

%% answer: 
\jdanswer{
$4$，$-2$
}
%% memo:
\memoanswer{
本题原卷及材料中未提供详解，待补充。
}

\item
%% number: 13
%% paperName: 1998年上海高考第 13 题 4 分
%% typeId: 3
%% type: 填空题
%% chapter: 磁场
%% point: 安培力
%% method: 无
%% score: 4
%% degree: 650
%% duplicateId: 0
%% body:
在同一平面上有 a、b、c 三根等间距平行放置的长直导线，依次载有电流强度为 $1\UA$、$2\UA$ 和 $3\UA$ 的电流，各电流的方向如图所示。则导线 a 所受的合力方向向 \tkanswer[1.2cm]{左}，导线 b 所受的合力方向向 \tkanswer[1.2cm]{右}。

\onepicture[h=4cm, align=right]{\includesvg{figs/13}}

%% answer: 
\jdanswer{
左，右
}
%% memo:
\memoanswer{
本题原卷及材料中未提供详解，待补充。
}

\item
%% number: 14
%% paperName: 1998年上海高考第 14 题 4 分
%% typeId: 3
%% type: 填空题
%% chapter: 交流电
%% point: 变压器
%% method: 无
%% score: 4
%% degree: 650
%% duplicateId: 0
%% body:
理想变压器原副线圈匝数比为 $N_{1}:N_{2} = 2:1$，原线圈接 $200\UV$ 交流电源、副线圈接额定功率为 $20\UW$ 的灯泡 L，灯泡正常发光。当电源电压降为 $180\UV$ 时，灯泡实际消耗功率与其额定功率之比为 \tkanswer[1.2cm]{0.81}，此时灯泡中的电流为 \tkanswer[1.6cm]{$0.18\UA$}（设灯泡电阻恒定）。

\onepicture[w=5.5cm, align=right]{\includesvg{figs/14}}

%% answer: 
\jdanswer{
0.81，$0.18\UA$
}
%% memo:
\memoanswer{
本题原卷及材料中未提供详解，待补充。
}

\item
%% number: 15
%% paperName: 1998年上海高考第 15 题 4 分
%% typeId: 3
%% type: 填空题
%% chapter: 功和能
%% point: 功率
%% method: 无
%% score: 4
%% degree: 450
%% duplicateId: 0
%% body:
某商场安装了一台倾角为 $30^{\circ}$ 的自动扶梯，该扶梯在电压为 $380\UV$ 的电动机带动下以 $0.4\Ums$ 的恒定速率向斜上方移动，电动机的最大输出功率为 $4.9\UkW$。不载人时测得电动机中的电流为 $5\UA$，若载人时扶梯的速率和不载人时相同，则这台自动扶梯可同时乘载的最多人数为 \tkanswer[1.2cm]{25}（设人的平均质量为 $60\Ukg$）。

%% answer: 
\jdanswer{
25
}
%% memo:
\memoanswer{
本题原卷及材料中未提供详解，待补充。
}

\item
%% number: 16
%% paperName: 1998年上海高考第 16 题 4 分
%% typeId: 3
%% type: 填空题
%% chapter: 曲线运动
%% point: 圆周运动的应用
%% method: 无
%% score: 4
%% degree: 450
%% duplicateId: 0
%% body:
质量为 $m$、带电量为 $+q$ 的小球用一绝缘细线悬于 O 点。开始时它在 A、B 之间来回摆动，OA、OB 与竖直方向 OC 的夹角均为 $\theta$（如图所示）。
\begin{enumerate}
	\item
	如果当它摆动到 B 点时突然施加一竖直向上的、大小为 $E = \frac{mg}{q}$ 的匀强电场，则此时线中拉力 $T_{1} =$ \tkanswer[2.2cm]{$0$}；
	\item
	如果这一电场是在小球从 A 点摆到最低点 C 时突然加上去的，则当小球运动到 B 点时线中的拉力 $T_{2} =$ \tkanswer[3.4cm]{$2mg(1 - \cos\theta)$}。
\end{enumerate}

\onepicture[h=5cm, align=right]{\includesvg{figs/16}}

%% answer: 
\jdanswer{
\begin{enumerate}
	\item
	$0$

	\item
	$2mg(1 - \cos\theta)$

\end{enumerate}
}
%% memo:
\memoanswer{
本题原卷及材料中未提供详解，待补充。
}

\item
%% number: 17
%% paperName: 1998年上海高考第 17 题 4 分
%% typeId: 3
%% type: 填空题
%% chapter: 功和能
%% point: 功率
%% method: 建立模型
%% score: 4
%% degree: 450
%% duplicateId: 0
%% body:
人的心脏每跳一次大约输送 $8 \times 10^{-5}\Umc$ 的血液，正常人血压（可看作心脏压送血液的压强）的平均值约为 $1.5 \times 10^{4}\UPa$，心跳约每分钟 70 次。据此估测心脏工作的平均功率约为 \tkanswer[1.9cm]{$1.4\UW$}。

%% answer: 
\jdanswer{
$1.4\UW$
}
%% memo:
\memoanswer{
本题原卷及材料中未提供详解，待补充。
}

\item
%% number: 18
%% paperName: 1998年上海高考第 18 题 6 分
%% typeId: 4
%% type: 实验题
%% chapter: 曲线运动
%% point: 研究平抛运动实验
%% method: 无
%% score: 6
%% degree: 650
%% duplicateId: 0
%% body:
在做“研究平抛运动”的实验时，让小球多次沿同一轨道运动，通过描点法画小球作平抛运动的轨迹。为了能较准确地描绘运动轨迹，下面列出了一些操作要求，将你认为正确的选项前面的字母填在横线上：\tkanswer[1.6cm]{ace}。

（a）通过调节使斜槽的末端保持水平

（b）每次释放小球的位置必须不同

（c）每次必须由静止释放小球

（d）记录小球位置用的木条（或凹槽）每次必须严格地等距离下降

（e）小球运动时不应与木板上的白纸（或方格纸）相接触

（f）将球的位置记录在纸上后，取下纸，用直尺将点连成折线

%% answer: 
\jdanswer{
ace
}
%% memo:
\memoanswer{
本题原卷及材料中未提供详解，待补充。
}

\item
%% number: 19
%% paperName: 1998年上海高考第 19 题 5 分
%% typeId: 4
%% type: 实验题
%% chapter: 电路
%% point: 多用表的使用
%% method: 无
%% score: 5
%% degree: 750
%% duplicateId: 0
%% body:
下图表示用多用表测电路中电流的实验，图中多用表测定的是 \tkanswer[1.8cm]{乙电阻}（填“甲电阻”、“乙电阻”或“总”）的电流，测得电流的大小是 \tkanswer[1.8cm]{$50\UmA$}。

\onepicture[w=8cm, align=right]{\includesvg{figs/19}}

%% answer: 
\jdanswer{
乙电阻，$50\UmA$
}
%% memo:
\memoanswer{
本题原卷及材料中未提供详解，待补充。
}

\item
%% number: 20
%% paperName: 1998年上海高考第 20 题 5 分
%% typeId: 4
%% type: 实验题
%% chapter: 机械波
%% point: 干涉衍射
%% method: 无
%% score: 5
%% degree: 550
%% duplicateId: 0
%% body:
如图是观察水面波衍射的实验装置，AC 和 BD 是两块挡板，AB 是一个孔，O 是波源，图中已画出波源所在区域波的传播情况，每两条相邻波纹（图中曲线）之间距离表示一个波长，则波经过孔之后的传播情况，下列描述中正确的是\xzanswer{ABC}

\onepicture[w=5cm]{\includesvg{figs/20}}

\fourchoices[answer=ABC]
{此时能明显观察到波的衍射现象}
{挡板前后波纹间距离相等}
{如果将孔 AB 扩大，有可能观察不到明显的衍射现象}
{如果孔的大小不变，使波源频率增大，能更明显观察到衍射现象}

%% answer: ABC
\jdanswer{
ABC
}
%% memo:
\memoanswer{
本题原卷及材料中未提供详解，待补充。
}

\item
%% number: 21
%% paperName: 1998年上海高考第 21 题 6 分
%% typeId: 4
%% type: 实验题
%% chapter: 气体性质
%% point: 验证玻意耳定律
%% method: 无
%% score: 6
%% degree: 650
%% duplicateId: 0
%% body:
某同学用一个水平固定着的、活塞横截面积已知的注射器和一个弹簧秤，测定当时的大气压。下面提供了四项必要的实验操作，请按正确的顺序将各项操作前的字母填在横线上：\tkanswer[1.6cm]{badc}。

（a）读出初始状态气体的体积 $V_{1}$

（b）将注射器活塞移到某一适当位置，用橡皮帽将注射器出口封住

（c）读出活塞被拉出一定距离后气体的体积 $V_{2}$ 和弹簧秤的拉力 $F$

（d）用弹簧秤将活塞拉出一定的距离

如果活塞的横截面积为 $S$，则测得大气压的表达式为 $p_{0} =$ \tkanswer[2.8cm]{$\frac{V_{2}F}{(V_{2} - V_{1})S}$}。

%% answer: 
\jdanswer{
badc，$\frac{V_{2}F}{(V_{2} - V_{1})S}$
}
%% memo:
\memoanswer{
本题原卷及材料中未提供详解，待补充。
}

\item
%% number: 22
%% paperName: 1998年上海高考第 22 题 7 分
%% typeId: 4
%% type: 实验题
%% chapter: 电路
%% point: 电路实验
%% method: 无
%% score: 7
%% degree: 550
%% duplicateId: 0
%% body:
如图所示，有 A、B、C 三个接线柱，A、B 之间接有内阻不计的 $10\UV$ 电源，手头有四个阻值完全相同的电阻，将它们适当组合，接在 A、C 和 C、B 之间，构成一个电路，使 A、C 间电压为 $6\UV$，C、B 间电压为 $4\UV$，试设计两种方案，分别画在图（a）、（b）中。

\twopicture[labelA=1998上海22a, labelB=1998上海22b, numA=（a）, numB=（b）, w=5cm]
{figs/22a.pdf}
{figs/22b.pdf}

%% answer: 
\jdanswer{
两方案的电路如图所示。

\onepicture[w=9cm]{\includesvg{figs/22_an}}
}
%% memo:
\memoanswer{
本题原卷及材料中未提供详解，待补充。
}

\item
%% number: 23
%% paperName: 1998年上海高考第 23 题 10 分
%% typeId: 6
%% type: 计算题
%% chapter: 光学
%% point: 光的折射
%% method: 无
%% score: 10
%% degree: 650
%% duplicateId: 0
%% body:
半径为 $R$ 的玻璃半圆柱体，横截面如图所示，圆心为 O。两条平行单色红光沿截面射向圆柱面，方向与底面垂直。光线 1 的入射点 A 为圆柱面的顶点，光线 2 的入射点 B，$\angle AOB = 60^{\circ}$。已知该玻璃对红光的折射率 $n = \sqrt{3}$。
\begin{enumerate}
	\item
	求两条光线经柱面和底面折射后的交点与 O 点的距离 $d$。
	\item
	若入射的是单色蓝光，则距离 $d$ 将比上面求得的结果大还是小？
\end{enumerate}

\onepicture[w=5cm, align=right]{\includesvg{figs/23}}

%% answer: 
\jdanswer{
\begin{enumerate}
	\item
	$d = \frac{R}{3}$

	\item
	$d$ 比上面结果小

\end{enumerate}
}
%% memo:
\memoanswer{
\onepicture[w=4.2cm, align=right]{\includesvg{figs/23_an}}

（1）光线 1 不偏折。

光线 2：$i = 60^{\circ}$。

$\sin r = \frac{\sin i}{n} = \frac{\frac{\sqrt{3}}{2}}{\sqrt{3}} = \frac{1}{2}$，解得 $r = 30^{\circ}$。

$i' = 60^{\circ} - r = 30^{\circ}$。

$\sin r' = n\sin i' = \frac{\sqrt{3}}{2}$，解得 $r' = 60^{\circ}$。

$OC = \frac{\frac{R}{2}}{\cos 30^{\circ}} = \frac{1}{\sqrt{3}}R$。

$d = OD = OC\tan 30^{\circ} = \frac{1}{\sqrt{3}}R \cdot \frac{1}{\sqrt{3}} = \frac{R}{3}$。

（2）$d$ 比上面结果小。

本题共 10 分。

（1）8 分。用到光线 1 不偏折 1 分，得出 $r = 30^{\circ}$ 2 分，得出 $r' = 60^{\circ}$ 2 分，得出 $d = \frac{R}{3}$ 3 分。

（2）2 分。正确答出 $d$ 变小的 2 分。
}

\item
%% number: 24
%% paperName: 1998年上海高考第 24 题 12 分
%% typeId: 6
%% type: 计算题
%% chapter: 气体性质
%% point: 气体连接体
%% method: 无
%% score: 12
%% degree: 450
%% duplicateId: 0
%% body:
如图所示，一个具有均匀横截面的不导热的封闭容器，被一不导热活塞分成 A、B 两部分。A、B 中充有同种理想气体，活塞可无摩擦地左右移动。开始时 A、B 的体积分别为 $V_{A} = 2V$，$V_{B} = V$，温度分别为 $T_{A}$ 和 $T_{B}$，两边压强均为 $p$，活塞处于平衡状态。现用某种方法使活塞能导热而发生移动，最后，两部分气体温度相同，两边的压强仍为 $p$。试求：
\begin{enumerate}
	\item
	最终状态时，A、B 两部分气体体积之比 $\frac{V_{A}'}{V_{B}'}$；
	\item
	最终状态时，A、B 两部分气体的温度 $T'$。
\end{enumerate}

\onepicture[w=6cm, align=right]{\includesvg{figs/24}}

%% answer: 
\jdanswer{
\begin{enumerate}
	\item
	$\frac{2T_{B}}{T_{A}}$

	\item
	$T' = \frac{3T_{A}T_{B}}{T_{A} + 2T_{B}}$

\end{enumerate}
}
%% memo:
\memoanswer{
（1）

$\frac{pV_{\mathrm{A}}}{T_{\mathrm{A}}} = \frac{pV_{\mathrm{A}}'}{T'}$    ①

$\frac{pV_{\mathrm{B}}}{T_{\mathrm{B}}} = \frac{pV_{\mathrm{B}}'}{T'}$    ②

①/②，并注意 $V_{A} = 2V$，$V_{B} = V$，得

$\frac{V_{A}'}{V_{B}'} = \frac{V_{A}T_{B}}{V_{B}T_{A}} = \frac{2T_{B}}{T_{A}}$    ③

（2）① + ②，

$\frac{pV_{\mathrm{A}}}{T_{\mathrm{A}}} + \frac{pV_{\mathrm{B}}}{T_{\mathrm{B}}} = \frac{p(V_{\mathrm{A}}' + V_{\mathrm{B}}')}{T'}$    ④

并注意 $V_{A} = 2V$，$V_{B} = V$ 及 $V_{A}' + V_{B}' = V_{A} + V_{B} = 3V$    ⑤

得 $T' = \frac{V_{A}' + V_{B}'}{\frac{V_{A}}{T_{A}} + \frac{V_{B}}{T_{B}}} = \frac{3V}{\frac{2V}{T_{A}} + \frac{V}{T_{B}}} = \frac{3T_{A}T_{B}}{T_{A} + 2T_{B}}$    ⑥

本题共 12 分。

（1）5 分。写出①、②式各 1 分，正确得出结果③ 3 分。

（2）7 分。写出④式 2 分，写出或用到⑤ 1 分，正确得出结果⑥式 4 分。
}

\item
%% number: 25
%% paperName: 1998年上海高考第 25 题 12 分
%% typeId: 6
%% type: 计算题
%% chapter: 电磁感应
%% point: 电磁感应中的力矩平衡
%% method: 无
%% score: 12
%% degree: 450
%% duplicateId: 0
%% body:
将一个矩形金属线框折成直角框架 abcdef，置于倾角为 $\alpha = 37^{\circ}$ 的斜面上，ab 边与斜面底线 MN 平行，如图所示，$ab = bc = cd = de = ef = fa = 0.2\Um$，线框总电阻为 $R = 0.02\UO$，ab 边和 de 边的质量均为 $m = 0.01\Ukg$，其余四边的质量忽略不计。框架可绕过 cf 点的固定轴转动。现从 $t = 0$ 时刻开始沿斜面向上加一随时间均匀增加的，范围足够大的匀强磁场，磁感强度与时间的关系 $B = 0.5t$（$\UT$），磁场方向与 cdef 面垂直。
\begin{enumerate}
	\item
	求线框中感应电流的大小，并在 ab 段导线上画出感应电流的方向。
	\item
	$t$ 为何值时框架将开始绕其固定轴转动？
\end{enumerate}

\onepicture[w=7cm, align=right]{\includesvg{figs/25}}

%% answer: 
\jdanswer{
\begin{enumerate}
	\item
	$1\UA$，a$\to$b

	\item
	$1.4\Us$

\end{enumerate}
}
%% memo:
\memoanswer{
（1）$E = \frac{\Delta \Phi}{\Delta t} = \frac{\Delta B}{\Delta t} \cdot cd \cdot de = 0.5 \times 0.2 \times 0.2\UV = 0.02\UV$

$I = \frac{E}{R} = \frac{0.02}{0.02}\UA = 1.0\UA$，方向从 a 到 b

（2）$f_{ab} = I \cdot ab \cdot B = 1.0 \times 0.2 \times 0.5t = 0.1t$，方向垂直斜面向上

$M_{f} = f_{ab} \cdot bc = 0.1t \times 0.2 = 0.02t$

$M_{G} = mg \cdot bc \cdot \cos\alpha + mg \cdot cd \cdot \sin\alpha = mg \cdot bc(\cos\alpha + \sin\alpha) = 0.01 \times 10 \times 0.2 \times (0.8 + 0.6)\UNcdotm = 0.028\UNcdotm$

令 $M_{f} = M_{G}$，$0.02t = 0.028$，$t = 1.4\Us$

本题共 12 分。

（1）4 分。正确得出 $I$ 值，3 分，正确给出方向 1 分。

（2）8 分。正确写出 $M_{f}$ 表达式 2 分，正确写出 $M_{G}$ 表达式 3 分，正确得出 $t$ 结果 3 分。
}

\item
%% number: 26
%% paperName: 1998年上海高考第 26 题 14 分
%% typeId: 6
%% type: 计算题
%% chapter: 运动和力
%% point: 动量和能量
%% method: 无
%% score: 14
%% degree: 250
%% duplicateId: 0
%% body:
用质量为 $M$ 的铁锤沿水平方向将质量为 $m$、长为 $l$ 的铁钉敲入木板，铁锤每次以相同的速度 $v_{0}$ 击钉，随即与钉一起运动并使钉进入木板一定距离。在每次受击进入木板的过程中，钉所受到的平均阻力为前一次受击进入木板过程所受平均阻力的 $k$ 倍（$k > 1$）。
\begin{enumerate}
	\item
	若敲击三次后钉恰好全部进入木板，求第一次进入木板过程中钉所受到的平均阻力。
	\item
	若第一次敲击使钉进入木板深度为 $l_{1}$，问至少敲击多少次才能将钉全部敲入木板？并就你的解答讨论要将钉全部敲入木板，$l_{1}$ 必须满足的条件。
\end{enumerate}

%% answer: 
\jdanswer{
\begin{enumerate}
	\item
	$f = \frac{M^{2}v_{0}^{2}(1 + k + k^{2})}{2(M + m)lk^{2}}$

	\item
	$n = \frac{\lg\left[1 - \frac{l}{l_{1}}\left(1 - \frac{1}{k}\right)\right]}{\lg\frac{1}{k}}$，若上式右边不是整数，$n$ 应为其取整加 1，若恰为整数，则不加 1；要使钉能全部钉入木板，应有 $l_{1} > \left(1 - \frac{1}{k}\right)l$。

\end{enumerate}
}
%% memo:
\memoanswer{
（1）$Mv_{0} = (m + M)v$    ①

$E_{k} = \frac{1}{2}(M + m)v^{2} = \frac{1}{2} \cdot \frac{M^{2}v_{0}^{2}}{M + m}$    ②

$E_{k} = fl_{1} = kfl_{2} = k^{2}fl_{3}$    ③

$l_{2} = \frac{1}{k}l_{1}$；$l_{3} = \frac{1}{k^{2}}l_{1}$

$l_{1} + l_{2} + l_{3} = l$    ④

得 $f = \frac{E_{k}}{l_{1}} = \frac{1}{2} \cdot \frac{M^{2}v_{0}^{2}}{M + m} \cdot \frac{1 + \frac{1}{k} + \frac{1}{k^{2}}}{l} = \frac{M^{2}v_{0}^{2}(1 + k + k^{2})}{2(M + m)lk^{2}}$    ⑤

（2）设敲 $n$ 次，钉子全部进入木板。

$E_{k} = fl_{1} = kfl_{2} = \cdots = k^{n - 1}fl_{n}$

$l_{1}\left(1 + \frac{1}{k} + \frac{1}{k^{2}} + \cdots + \frac{1}{k^{n - 1}}\right) = l$    ⑥

$\frac{l}{l_{1}} = 1 + \frac{1}{k} + \frac{1}{k^{2}} + \cdots + \frac{1}{k^{n - 1}} = \frac{1 - 1/k^{n}}{1 - 1/k}$    ⑦

得 $n = \frac{\lg\left[1 - \frac{l}{l_{1}}\left(1 - \frac{1}{k}\right)\right]}{\lg\frac{1}{k}}$    ⑧

若上式右边不是整数，$n$ 应为其取整加 1，若恰为整数，则不加 1。

⑦式右边随 $n$ 的增大而增大，但总是小于 $\frac{1}{1 - 1/k}$，因而当 $l_{1}$ 太小时，无论多大的 $n$ 也不能使⑦式成立，故要使钉能全部钉入木板，应有

$\frac{l}{l_{1}} < \frac{1}{1 - 1/k}$，即 $l_{1} > \left(1 - \frac{1}{k}\right)l$    ⑨

本题共 14 分。

（1）7 分。写出①、②、③式各 1 分，写出或用到④ 1 分，得出结果⑤ 3 分。

（2）7 分。得出⑥式 1 分，得出⑦式 2 分，得出⑧式 2 分（未写 $n$ 取整不扣分），得出⑨式 2 分。
}

\end{enumerate}

\end{document}
